% Proposed replacement abstract (2026-09-15). Not yet wired into main.tex.
% Rationale: leads with the practitioner's question, states the finding with the
% numbers the tables actually support (within-0.05 at 30 samples rather than
% "saturate within a few dozen"), and promotes the three less obvious results
% (detailed prompt = an order of magnitude of data; red-teaming levels rather
% than adds; generator ranking does not transfer) from footnotes to claims.
\begin{abstract}
Linear activation probes are a cheap way to monitor deployed language models,
and their training data is almost always synthetic. Two questions decide the
cost of every such monitor: how many synthetic examples does the probe need,
and how should they be written? We answer both across three monitoring
concepts---\emph{high-stakes} situations, assistant replies \emph{harmful} to a
person, and replies that do not follow the user's \emph{instruction}---while
holding the probe model, layer, architecture, generators, prompt structure, and
evaluation protocol fixed, on fourteen held-out evaluation distributions. We
measure learning curves over 10--590 synthetic samples for two data sources,
direct prompting and an adversarial red-teaming scaffold that retrains on the
probe's own mistakes, each under a \emph{general} prompt that states only the
concept and a \emph{detailed} prompt that also describes the deployment
distribution. The concept, not the recipe, sets the shape of the curve: with
30 samples the high-stakes and harmfulness probes are within 0.05 AUROC of
their 590-sample level, whereas the instruction-following probe is 0.21 below
it, is not learnable at all from general prompts, and needs several hundred
detailed-prompt samples; a split-level analysis traces this to the number of
distinct failure modes the concept spans. The recipe sets the level: a
detailed prompt is worth an order of magnitude of data on the splits it
describes, with gains ordered exactly by the headroom the general prompt
leaves (+0.14 to +0.18 on instruction following, $\le$+0.03 on harmfulness);
adversarially found data adds up to +0.10 on the hard concept but only
levels probes toward a fixed band on the easy ones; and which generator
writes the best data reverses across concepts and budgets. We release the
evaluation suites, dev sets, and generated training sets, and give a
practical rule: measure the curve on held-out distributions before fixing a
budget, because the concept decides where it bends.
\end{abstract}

% ---------------------------------------------------------------------------
% Proposed replacement abstract v2 (2026-09-17). Not wired into main.tex.
% Rationale: leads with the one non-obvious claim (recipe moves the level,
% distribution sets the bend); ~210 words instead of 330; drops the
% "0.08--0.21 below at 30 samples" figure, which is the 80-sample cut, not 30;
% quotes the LLM-prompt gain against the hand-written description (+0.14 on
% Ant-HH) instead of against the 50-sample probe (+0.29); states the
% headroom-normalised ordering, which needs one sentence added to Sec. 5.1
% (fraction of the 10->590 gain reached by 80 samples: hs 0.85-0.90,
% harm 0.74-0.96, instr 0.24-0.70, computed from Table sizecurves).
\begin{abstract}
Activation probes that monitor deployed language models are trained almost
entirely on synthetic data, so two questions set the cost of every monitor:
how many samples, and how to write them. We measure learning curves over
10--590 synthetic samples for three monitoring concepts---\hs{} situations,
replies \harm{} to a person, and replies that do not follow the user's
\instr{}---on fourteen held-out distributions and four probe models from 1B to
27B parameters, under direct prompting and a red-teaming scaffold, with prompts
that know nothing about the deployment traffic and prompts an LLM writes from
samples of it. The recipe moves the level of the curve; the distribution
decides where it bends. Cutting 590 samples to 80 costs \hs{} probes at most
0.02 AUROC and \harm{} probes at most 0.06 but \instr{} probes 0.08--0.21, on
every probe model and after normalising out headroom. Within one concept,
saturation sizes differ by an order of magnitude between distributions; five
deliberately different prompts written from one distribution's samples differ
in final AUROC by up to 0.23 yet saturate at similar sizes; and within one
distribution, subpopulations of samples settle at 60, 120, 300 samples, or
not at all, replicably, with no surface feature of the text predicting which.
Red-teaming levels the easy concepts toward a fixed band and adds only to the
hard one. Data needs must be measured on the deployment distribution, and we
release the evaluation suites, dev sets, and generated sets to do so.
\end{abstract}
